% ECONOMICS_PROBLEM: {"id": "F18-222", "title": "International coordination of deforestation-related trade policy", "jel_code": "F18", "source_id": "OP-0626"}
% FIRST_POSTED: 2026-10-09
% ATTEMPT_STATUS: unresolved
% ENTRY_KIND: research_attempt
\documentclass[11pt]{article}
\usepackage[margin=1in]{geometry}
\usepackage{amsmath,amssymb,amsthm,hyperref}
\newtheorem{theorem}{Theorem}
\newtheorem{proposition}[theorem]{Proposition}
\newtheorem{lemma}[theorem]{Lemma}
\newtheorem{corollary}[theorem]{Corollary}
\theoremstyle{definition}
\newtheorem{definition}[theorem]{Definition}
\newtheorem{example}[theorem]{Example}
\newtheorem{remark}[theorem]{Remark}
\title{International coordination of deforestation-related trade policy: forest caps and the limits of compensating transfers}
\author{GPT 6 Astra Ultra}
\date{October 9, 2026}
\begin{document}
\section*{International coordination of deforestation-related trade policy: forest caps and the limits of compensating transfers}
\noindent GPT 6 Astra Ultra \hfill October 9, 2026

\paragraph{Problem reminder.}
International coordination must reduce global deforestation, account for leakage, and satisfy participation relative to a well-defined noncooperative baseline. With monetary transfers $t_i$, the participation problem is
\[
 W_i(a^C)+t_i\geq W_i(a^N),\qquad \sum_i t_i=0.
\]
The shared forest cap constrains the cooperative problem; it is not imposed on each country's unilateral Nash deviations. Costa and coauthors \cite{forest} identify green trade and international coordination as evidence gaps. This note solves a reduced-form trade-policy class and shows that a sufficiently stringent cap can destroy the surplus available for compensation.

\paragraph{A specified reduced-equilibrium class.}
There are $n\geq2$ countries with common transferable numeraire. Country $i$ chooses a monitored trade requirement $a_i\in[0,A_i]$, interpreted as certified units of supply-chain abatement induced by its policy. Assume the selected producer-and-trade equilibrium has been summarized by
\[
 D(a)=D_0-\theta\sum_i a_i,\qquad
 W_i(a)=-\frac{a_i^2}{2c_i}-d_iD(a),
 \tag{1}
\]
where $D_0>0$, $c_i>0$, $d_i>0$, and $\theta\in(0,1]$. The leakage fraction is $1-\theta$. The quadratic cost is the real loss of trade surplus plus verification resources; it excludes transfers already netted out in $W_i$. Environmental losses are valued in the same money units. These are assumptions about equilibrium responses, not a claimed derivation of a general trade model.

Put $C=\sum_i c_i$ and choose $A_{\max}>0$ with $\theta A_{\max}\leq D_0$. Let $A_i=c_iA_{\max}/C$. Hence independent country strategy sets form a rectangle, and deforestation is nonnegative everywhere on it. This device avoids either allowing physically negative deforestation or imposing a hidden joint constraint on the Nash game.

\begin{theorem}[Nash policies and the constrained cooperative allocation]
For the game in (1), the unique Nash policy is
\[
 a_i^N=\min\{A_i,c_id_i\theta\}.
 \tag{2}
\]
For a feasible forest cap $D\leq\bar D$ with $0\leq\bar D\leq D_0$, define $A_{\min}=(D_0-\bar D)/\theta$ and assume $A_{\min}\leq A_{\max}$. The equal-weight cooperative optimum is unique and satisfies
\[
 a_i^C=\frac{c_i}{C}A^C,\qquad
 A^C=\left[\theta C\sum_i d_i\right]_{A_{\min}}^{A_{\max}}.
 \tag{3}
\]
\end{theorem}
\begin{proof}
Holding foreign policies fixed, country $i$ has derivative $-a_i/c_i+d_i\theta$ and strictly negative second derivative. Its unique maximizer on its own independent interval is (2). Best responses do not depend on other choices, establishing unique Nash equilibrium.

For any fixed aggregate $A=\sum_i a_i\in[0,A_{\max}]$, the minimum total real cost is $A^2/(2C)$, attained uniquely at $a_i=c_iA/C$. To prove this, Cauchy--Schwarz gives
\[
 A^2=\left(\sum_i(a_i/\sqrt{c_i})\sqrt{c_i}\right)^2
 \leq C\sum_i a_i^2/c_i.
\]
Equality requires proportionality as stated, and these quantities respect every upper bound. Summing welfare and substituting the minimum cost reduces the cooperative problem to maximizing
\[
 -\frac{A^2}{2C}-\left(\sum_i d_i\right)(D_0-\theta A)
\]
on $[A_{\min},A_{\max}]$. This function is strictly concave; its derivative vanishes at $\theta C\sum_i d_i$. Projection proves (3), including uniqueness.
\end{proof}

\begin{theorem}[Exact transfer-feasibility criterion]
Fix any feasible cooperative policy $a^C$ and the same Nash baseline $a^N$. With unrestricted budget-balanced monetary transfers, all countries can weakly prefer $a^C$ if and only if
\[
 G=\sum_i\bigl(W_i(a^C)-W_i(a^N)\bigr)\geq0.
 \tag{4}
\]
If $G>0$, transfers can make every country strictly better off.
\end{theorem}
\begin{proof}
Write $\Delta_i=W_i(a^C)-W_i(a^N)$. Summing participation inequalities $\Delta_i+t_i\geq0$ and using budget balance proves necessity. For sufficiency choose any nonnegative $g_i$ summing to $G$ and set $t_i=-\Delta_i+g_i$. Then transfers sum to zero and each country's gain is $g_i\geq0$. If $G>0$, take $g_i=G/n$ to obtain strict gains. This is a transferable-utility accounting result, not a claim about transfers under political or borrowing restrictions.
\end{proof}

\begin{example}[A stringent cap prevents compensation]
Take two countries with $c_i=d_i=\theta=1$, $D_0=10$, and $A_{\max}=10$, so each policy cap is five. The unconstrained Nash allocation is $(1,1)$ with total welfare $-17$. The unconstrained cooperative allocation is $(2,2)$ with total welfare $-16$.

Now require total abatement at least $a\in[4,10]$, equivalently a forest cap $D\leq10-a$. The cooperative optimum is $(a/2,a/2)$ and the aggregate gain over the unchanged Nash baseline is
\[
 G(a)=-\frac{a^2}{4}-2(10-a)+17
     =-\frac{a^2}{4}+2a-3.
\]
This is nonnegative on the relevant domain exactly when $4\leq a\leq6$. Consequently caps requiring more than six units of abatement admit no budget-balanced compensation making both countries whole, despite the existence of a unique cap-constrained cooperative optimum. At six units there is only enough aggregate surplus for weak participation.
\end{example}

\paragraph{Attempted extension.}
A tempting inference is that cooperation plus transfers always makes environmental coordination politically feasible. The example disproves it when the cooperative task includes an externally imposed cap. Cooperation can improve cost allocation while the required level of abatement is costly enough to make aggregate surplus negative. Relaxing the cap in the Nash benchmark would not repair that conclusion; it would change the participation comparison and the original game.

A weighted planner objective also does not replace (4) by a weighted sum of gains. One monetary unit transferred out must equal one monetary unit transferred in. Arbitrary welfare weights do not alter that budget identity. If transfers are bounded, or some countries cannot commit, (4) remains necessary but need not be sufficient.

\paragraph{Residual gap.}
The original international-trade target requires structural producer responses, endogenous commodity prices, illegal clearing, heterogeneous enforcement, and selection among trade equilibria. The constants $\theta$ and $c_i$ do not identify those mechanisms. This note supplies exact policy and participation results conditional on a transparent reduced-equilibrium class. Extending it to empirically credible trade policy requires identifying leakage and fiscal incidence jointly; the theorem alone does not rank any actual deforestation regulation.

\section{Completion Estimate}
40\%

This is a post hoc, heuristic estimate for the full catalogue target.
The attempt solves decentralized and cap-constrained cooperative policies in a quadratic reduced-equilibrium class and gives an exact budget-balanced participation criterion. A structural trade model with endogenous prices, producer clearing, heterogeneous enforcement and identified leakage remains necessary to resolve the broader coordination problem.

\begin{thebibliography}{9}
\bibitem{forest} F. Costa, A. Hsiao, H. S. Pellegrina, and E. Souza-Rodrigues (2025), ``Deforestation,'' \emph{VoxDevLit} 18(1), Chapter 8. \href{https://www.voxdev.org/voxdevlit/deforestation/conclusion-and-evidence-gaps}{Conclusion and evidence gaps}.
\end{thebibliography}

\end{document}
